% =============================================================================
%  Sección 1: Introducción
%  Sustituye el contenido de \lipsum con tu propio texto.
% =============================================================================

\section{Introducción}

La \termino{densificación urbana} constituye uno de los ejes centrales del debate contemporáneo en urbanismo.
Frente al modelo de expansión horizontal que ha dominado el crecimiento metropolitano durante el siglo XX,
emergen nuevas aproximaciones orientadas a la compacidad, la mixticidad de usos y la movilidad no motorizada
como vectores de un desarrollo más equitativo y ambientalmente responsable \cite{lore2024densidad}.

\lipsum[2]

\subsection{Contexto metropolitano}

En el caso de la \zmvm, la presión demográfica combinada con la ausencia de instrumentos de planeación vinculantes
ha producido un patrón de crecimiento que \termino{fragmenta el territorio} y eleva exponencialmente los costos
de provisión de infraestructura y servicios urbanos \cite{lore2024zmvm}.

\lipsum[3]

\subsection{Objetivo del ensayo}

El presente ensayo tiene por objetivo evaluar tres modelos de densificación —el modelo nórdico de
\textit{transit-oriented development} (TOD), el modelo latinoamericano de intervención en barrios
consolidados, y el modelo de corredores urbanos— bajo los criterios de equidad, eficiencia energética
y viabilidad institucional en el contexto mexicano.

\begin{itemize}
    \item Analizar los fundamentos teóricos de cada modelo de densificación.
    \item Identificar experiencias internacionales exitosas y sus condiciones de transferibilidad.
    \item Proponer lineamientos de política pública aplicables a la \zmvm.
\end{itemize}

\lipsum[4]
